\section{Dependency injection}
\label{sec:dependency-injection}

The dependency-injection utilities provide a small service workflow for Unity projects. They do not require a container package: services are registered with \texttt{ServiceLocator}, fields or properties are marked with \texttt{[Inject]}, and \texttt{Injector} supplies matching services.

\subsection{Registering and resolving services}
\texttt{ServiceLocator} registers one object per type and exposes \texttt{TryGet}, \texttt{GetAllServices}, and \texttt{Clear}. Registration is explicit, so startup order remains visible in the scene.

\begin{lstlisting}
using MLGWorks.Utils.DependencyInjection;

ServiceLocator.Register(typeof(GameSettings), new GameSettings());

if (ServiceLocator.TryGet(typeof(GameSettings), out object value))
{
    GameSettings settings = (GameSettings)value;
}
\end{lstlisting}

\subsection{Injection}
Mark a field or property with \texttt{[Inject]} and call \texttt{Injector.Inject(target)}. \texttt{DIInitializer} performs this work for scene objects, including inactive objects, and offers \texttt{InitializeNow()} when a caller needs to trigger the pass explicitly.

\begin{lstlisting}
public sealed class HudController : MonoBehaviour, IInitializable
{
    [Inject] private GameSettings settings;

    public void Initialize()
    {
        ApplyTheme(settings.Theme);
    }
}
\end{lstlisting}

\subsection{Startup and auto-registration}
\texttt{ServiceBootstrapper} is intended for registering core services early in the Unity lifecycle. A MonoBehaviour marked with \texttt{[DIService]} can be discovered and registered by the initializer. Implement \texttt{IInitializable} when a service needs a callback after injection is complete.

The attribute types live under \texttt{Dependency Injection/Attributes}; the runtime types remain in the parent namespace \texttt{MLGWorks.Utils.DependencyInjection}.
